\begin{examplebox}
	\textbf{\textcolor{CExample}{例 1（基础）}}

	\textbf{\textcolor{CExample}{题目：}}
	已知 $\sin\alpha=\dfrac{3}{5}$，$\alpha\in\left(0,\dfrac{\pi}{2}\right)$，$\cos\beta=-\dfrac{5}{13}$，$\beta\in\left(\pi,\dfrac{3\pi}{2}\right)$，求 $\cos(\alpha-\beta)$ 的值。

	\textbf{\textcolor{CExample}{解题步骤：}}
	\begin{description}[style=nextline, leftmargin=2.8em, labelsep=0.5em, itemsep=0.45em]
		\item[\textbf{第一步（求 $\cos\alpha$）：}] 由 $\sin^2\alpha+\cos^2\alpha=1$，得 $\cos\alpha=\pm\sqrt{1-\sin^2\alpha}=\pm\dfrac{4}{5}$。因为 $\alpha\in\left(0,\dfrac{\pi}{2}\right)$，$\cos\alpha>0$，故 $\cos\alpha=\dfrac{4}{5}$。
			\par\textbf{理由：} 同角三角函数的平方关系，符号由象限决定。

		\item[\textbf{第二步（求 $\sin\beta$）：}] $\sin\beta=\pm\sqrt{1-\cos^2\beta}=\pm\sqrt{1-\left(-\dfrac{5}{13}\right)^2}=\pm\dfrac{12}{13}$。因为 $\beta\in\left(\pi,\dfrac{3\pi}{2}\right)$，第三象限正弦为负，所以 $\sin\beta=-\dfrac{12}{13}$。
			\par\textbf{理由：} 平方关系结合象限符号。

		\item[\textbf{第三步（代入余弦差公式）：}]
			\[
				\begin{aligned}
					\cos(\alpha-\beta) &= \cos\alpha\cos\beta+\sin\alpha\sin\beta \\
					&= \dfrac{4}{5}\times\left(-\dfrac{5}{13}\right)+\dfrac{3}{5}\times\left(-\dfrac{12}{13}\right) \\
					&= -\dfrac{20}{65}-\dfrac{36}{65} \\
					&= -\dfrac{56}{65}.
			\end{aligned}
			\]
			\par\textbf{理由：} 余弦差公式 $\cos(\alpha-\beta)=\cos\alpha\cos\beta+\sin\alpha\sin\beta$。
	\end{description}

	\textbf{\textcolor{CExample}{关键结论：}}
	\textbf{$\cos(\alpha-\beta)=-\dfrac{56}{65}$}.

	\medskip
	\textbf{\textcolor{CExample}{例 2（稍变形）}}

	\textbf{\textcolor{CExample}{题目：}}
	已知 $\tan(\alpha+\beta)=2$，$\tan\alpha=1$，求 $\tan\beta$ 的值。

	\textbf{\textcolor{CExample}{解题步骤：}}
	\begin{description}[style=nextline, leftmargin=2.8em, labelsep=0.5em, itemsep=0.45em]
		\item[\textbf{第一步（展开正切和公式）：}]
			\[
				\tan(\alpha+\beta)=\frac{\tan\alpha+\tan\beta}{1-\tan\alpha\tan\beta}.
			\]
			\par\textbf{理由：} 正切和角公式。

		\item[\textbf{第二步（代入已知值）：}] 代入 $\tan(\alpha+\beta)=2$，$\tan\alpha=1$，得
			\[
				2=\frac{1+\tan\beta}{1-1\cdot\tan\beta},
			\]
			即
			\[
				2=\frac{1+\tan\beta}{1-\tan\beta}.
			\]
			\par\textbf{理由：} 数值代入公式。

		\item[\textbf{第三步（解方程）：}] 解方程 $2=\dfrac{1+\tan\beta}{1-\tan\beta}$，两边乘以 $1-\tan\beta$ 得 $2(1-\tan\beta)=1+\tan\beta$，展开得 $2-2\tan\beta=1+\tan\beta$，移项得 $1=3\tan\beta$，所以 $\tan\beta=\dfrac{1}{3}$。
			\par\textbf{理由：} 代数运算求解。
	\end{description}

	\textbf{\textcolor{CExample}{关键结论：}}
	\textbf{$\tan\beta=\dfrac{1}{3}$}.

	\medskip
	\textbf{\textcolor{CExample}{例 3（综合应用）}}

	\textbf{\textcolor{CExample}{题目：}}
	在 $\triangle ABC$ 中，$A,B,C\neq\frac{\pi}{2}$，求证：$\tan A+\tan B+\tan C=\tan A\tan B\tan C$。

	\textbf{\textcolor{CExample}{解题步骤：}}
	\begin{description}[style=nextline, leftmargin=2.8em, labelsep=0.5em, itemsep=0.45em]
		\item[\textbf{第一步（利用内角和）：}] 在 $\triangle ABC$ 中，$A+B+C=\pi$，从而 $A+B=\pi-C$。两边取正切得 $\tan(A+B)=\tan(\pi-C)$。因为 $\tan(\pi-C)=-\tan C$（正切周期为 $\pi$，且 $\tan(\pi-\theta)=-\tan\theta$），所以 $\tan(A+B)=-\tan C$。
			\par\textbf{理由：} 三角恒等式与诱导公式。

		\item[\textbf{第二步（展开正切和）：}] 对左侧应用正切和公式：
			\[
				\tan(A+B)=\frac{\tan A+\tan B}{1-\tan A\tan B}.
			\]
			因此 $\dfrac{\tan A+\tan B}{1-\tan A\tan B}=-\tan C$。
			\par\textbf{理由：} 正切和角公式。

		\item[\textbf{第三步（变形整理）：}] 将上式变形：
			\[
				\begin{aligned}
					\tan A+\tan B &= -\tan C\,(1-\tan A\tan B) \\
					&= -\tan C+\tan A\tan B\tan C.
			\end{aligned}
			\]
			移项得 $\tan A+\tan B+\tan C=\tan A\tan B\tan C$。
			\par\textbf{理由：} 代数运算，将项移至同侧。
	\end{description}

	\textbf{\textcolor{CExample}{关键结论：}}
	\textbf{恒等式 $\tan A+\tan B+\tan C=\tan A\tan B\tan C$ 得证。}
\end{examplebox}
